\documentclass{article}\usepackage[]{graphicx}\usepackage[]{xcolor}
% maxwidth is the original width if it is less than linewidth
% otherwise use linewidth (to make sure the graphics do not exceed the margin)
\makeatletter
\def\maxwidth{ %
  \ifdim\Gin@nat@width>\linewidth
    \linewidth
  \else
    \Gin@nat@width
  \fi
}
\makeatother

\definecolor{fgcolor}{rgb}{0.345, 0.345, 0.345}
\newcommand{\hlnum}[1]{\textcolor[rgb]{0.686,0.059,0.569}{#1}}%
\newcommand{\hlsng}[1]{\textcolor[rgb]{0.192,0.494,0.8}{#1}}%
\newcommand{\hlcom}[1]{\textcolor[rgb]{0.678,0.584,0.686}{\textit{#1}}}%
\newcommand{\hlopt}[1]{\textcolor[rgb]{0,0,0}{#1}}%
\newcommand{\hldef}[1]{\textcolor[rgb]{0.345,0.345,0.345}{#1}}%
\newcommand{\hlkwa}[1]{\textcolor[rgb]{0.161,0.373,0.58}{\textbf{#1}}}%
\newcommand{\hlkwb}[1]{\textcolor[rgb]{0.69,0.353,0.396}{#1}}%
\newcommand{\hlkwc}[1]{\textcolor[rgb]{0.333,0.667,0.333}{#1}}%
\newcommand{\hlkwd}[1]{\textcolor[rgb]{0.737,0.353,0.396}{\textbf{#1}}}%
\let\hlipl\hlkwb

\usepackage{framed}
\makeatletter
\newenvironment{kframe}{%
 \def\at@end@of@kframe{}%
 \ifinner\ifhmode%
  \def\at@end@of@kframe{\end{minipage}}%
  \begin{minipage}{\columnwidth}%
 \fi\fi%
 \def\FrameCommand##1{\hskip\@totalleftmargin \hskip-\fboxsep
 \colorbox{shadecolor}{##1}\hskip-\fboxsep
     % There is no \\@totalrightmargin, so:
     \hskip-\linewidth \hskip-\@totalleftmargin \hskip\columnwidth}%
 \MakeFramed {\advance\hsize-\width
   \@totalleftmargin\z@ \linewidth\hsize
   \@setminipage}}%
 {\par\unskip\endMakeFramed%
 \at@end@of@kframe}
\makeatother

\definecolor{shadecolor}{rgb}{.97, .97, .97}
\definecolor{messagecolor}{rgb}{0, 0, 0}
\definecolor{warningcolor}{rgb}{1, 0, 1}
\definecolor{errorcolor}{rgb}{1, 0, 0}
\newenvironment{knitrout}{}{} % an empty environment to be redefined in TeX

\usepackage{alltt}
\usepackage{hyperref}
\usepackage[margin=1in]{geometry}
\usepackage{booktabs}
\IfFileExists{upquote.sty}{\usepackage{upquote}}{}
\begin{document}
% \SweaveOpts{concordance=TRUE}

\title{Supp material for Wildchrokie and CoringTreespotters}
\author{Christophe Rouleau-Desrochers}
\date{\today}
\maketitle



% <><><><><><><><><><><><><><><><><><><><><><><><><><><><><><><><><><><><><><><>
\section{Wildchrokie}
% <><><><><><><><><><><><><><><><><><><><><><><><><><><><><><><><><><><><><><><>

% <><><><><><><><><><><><><><><><><><><><><><><><><><><><><><><><><><><><><><><>
\section{CoringTreespotters}
% <><><><><><><><><><><><><><><><><><><><><><><><><><><><><><><><><><><><><><><>
\subsection{Age allometry}
It is commonplace in dendrochronoly to perform standardization, i.e. selecting and fitting a statistical model to the annual growth increments of each tree separately to remove the non climatic growth influences such as those associated with tree age (CITE Schofield 2016 and Fritts 1976). This filters our non climatic growth influences. However, this is most often performed on very long series, where age allometry is expected to affect the ring width as the tree's diameter gets wider. We did not detrend our ring chronologies because we had 19 trees that had decreasing ring widths as a function of tree age versus 25 that showed a positive trend, for our 9 years of interest (Figure \ref{fig:rwByAge}. For the whole chronologies, we found that 27 trees had a decreasing ring width trend with age compared to 17 that had a positive trend.

% Figure Ring width across years
\begin{figure}[!ht]
\includegraphics[width=1\textwidth]{/Users/christophe_rouleau-desrochers/github/coringtreespotters/analyses/figures/empiricalData/RWXage_bySpp9Yrs1page.pdf}
\caption{Ring width as a function of tree age at the corresponding year. The line represents the linear model ring width as a function of age. The dots are the empirical data and the mean age is the averaged estimated age of that tree during the studied period.}
\label{fig:rwByAge}
\end{figure}



% Climate lmer output


\end{document}
